
\medskip
\phantomsection\label{par:sga3-VIA-5.4}%
\noindent\textbf{5.4.---}%
\footnote{\label{note:sga3-VIA-53a}Editor's note: we have expanded
the original in what follows; in particular, we have added
\hyperref[prop:sga3-VIA-5.4.1]{Proposition~5.4.1}.}
Let \(u:G\to H\) be a quasi-compact morphism between \(A\)-groups
locally of finite type, such that the kernel \(N\) of \(u\) is flat
over \(A\). In this case, by \SGARefLink{thm:sga3-VIA-3.3.2}{3.3.2} and \SGARefLink{prop:sga3-VIB-5.2}{5.2}, the quotient \(A\)-group
\[
  C=N\backslash G
\]
exists, and the morphism \(p:G\to C\) is faithfully flat and locally
of finite presentation. On the other hand, by \SGARefLink{sga3:iv:remark:5:2:6}{Expos\'e~IV, 5.2.6},
\(u\) induces a monomorphism
\[
  v:C\longrightarrow H,
\]
which is quasi-compact (because \(u\) is quasi-compact and
\(G\to C\) is surjective; cf.\ EGA~IV\(_1\), 1.1.3), and is
therefore a closed immersion by \hyperref[prop:sga3-VIA-2.5.2]{2.5.2}. We have thus obtained the
following proposition.

\medskip
\phantomsection\label{prop:sga3-VIA-5.4.1}%
\noindent\textbf{Proposition 5.4.1.---}
\emph{Let \(u:G\to H\) be a quasi-compact morphism between
\(A\)-groups locally of finite type, such that
\(N=\ker u\) is flat over \(A\). Then \(u\) admits the factorization}
\begin{center}
  \phantomsection\label{diag:sga3-VIA-5.4.1-factorization}%
  \begin{tikzcd}[row sep=large,column sep=large]
    G \arrow[r,"u"] \arrow[d,"p"'] & H \\
    N\backslash G \arrow[ur,"i"']
  \end{tikzcd}
\end{center}
\emph{where \(p\) is faithfully flat and locally of finite
presentation, and \(i\) is a closed immersion.}
\par\medskip

Suppose in addition that \(G\) is flat over \(A\). Then, by \SGARefLink{thm:sga3-VIA-3.3.2}{3.3.2},
\(C=N\backslash G\) is flat over \(A\), and hence the quotient
\[
  X=C\backslash H
\]
exists in \((\mathbf{Sch}_{/A})\) and represents the quotient
\((\mathrm{fppf})\)-sheaf \(\widetilde C\backslash H\), while
\[
  q:H\longrightarrow X
\]
is a \(C\)-torsor. Consequently, if
\[
  e:\operatorname{Spec}A\longrightarrow H
\]
denotes the unit section of \(H\), then \(v\) induces an isomorphism
of \((\mathrm{fppf})\)-sheaves between \(\widetilde C\) and the fiber
product of \(q\) with
\[
  q\circ e:\operatorname{Spec}A\longrightarrow X,
\]
which is represented by a closed subscheme of \(H\). Consequently,
\(v\) is an isomorphism of \(C\) onto a closed group subscheme
\(K\) of \(H\), equal to the stabilizer of the \(A\)-point
\(q\circ e\) of \(X\).%
\begingroup
\renewcommand{\thefootnote}{\fnsymbol{footnote}}%
\footnote[1]{\label{note:sga3-VIA-5.4-K-ambient}Translator's note:
the French source declares \(e:\operatorname{Spec}A\to G\), although
the displayed composition \(q\circ e\) requires the unit section of
\(H\), and then prints ``a closed group subscheme \(K\) of \(G\).''
Here \(v:C\to H\), the displayed fiber product is represented inside
\(H\), and the next paragraph identifies \(K\) with
\(\ker(H\to\overline H)\); accordingly, \(e\) is taken to be the unit
section of \(H\), and the intended ambient group of \(K\) is \(H\).}%
\endgroup{}
(This also gives another proof that every quasi-compact monomorphism
\(v:C\to H\) between \(A\)-groups locally of finite type is a closed
immersion; cf.\ \SGARefLink{prop:sga3-VIA-2.5.2}{2.5.2} and Expos\'e~VI\(_B\), \SGARefLink{cor:sga3-VIB-1.4.2}{1.4.2}.)

Suppose further that \(C\) is a normal subgroup of \(H\). In this
case, the \(A\)-group
\[
  \overline H=C\backslash H
\]
is the cokernel, in the category of \(A\)-groups, of the morphism
\(u:G\to H\), and \(K\) is the kernel of the morphism
\[
  H\longrightarrow\overline H.
\]
When \(G\) and \(H\) are abelian \(A\)-groups, \(K\) is the image of
\(u\) in the category of abelian \(A\)-groups, whereas
\[
  C=(\ker u)\backslash G
\]
is the coimage of \(u\). Taking into account the isomorphism
\[
  C\xrightarrow{\ \sim\ }K
\]
just established, we obtain the following.%
\footnote[53]{\label{note:sga3-VIA-53b}Editor's note: we have added
the number \SGARefLink{thm:sga3-VIA-5.4.2}{5.4.2} to this theorem.}

\medskip
\phantomsection\label{thm:sga3-VIA-5.4.2}%
\noindent\textbf{Theorem 5.4.2.---}
\emph{Let \(k\) be a field. The category of commutative algebraic
\(k\)-groups is abelian.}

Indeed, when \(k\) is a field, \(\ker(u)\) is flat over \(k\),
whatever \(u\) may be.

\noindent
\footnote{\label{note:sga3-VIA-54}Editor's note: we have added what
follows.}
Note that the full subcategory of commutative affine algebraic
\(k\)-groups is thick. Indeed, consider an exact sequence of
commutative algebraic \(k\)-groups
\[
  1\longrightarrow N\longrightarrow G\longrightarrow
  G/N\longrightarrow 1.
\]
If \(G\) is affine, then \(N\) is clearly affine, and \(G/N\) is
also affine by a theorem of Chevalley; cf.\ Expos\'e~VI\(_B\),
\SGARefLink{thm:sga3-VIB-11.17}{11.17}. Conversely, if \(N\) and \(G/N\) are affine, then \(G\) is
also affine by Expos\'e~VI\(_B\), \SGARefLink{concl:sga3-VIB-9.2-viii}{9.2(viii)}. We therefore obtain the
following corollary.

\medskip
\phantomsection\label{cor:sga3-VIA-5.4.3}%
\noindent\textbf{Corollary 5.4.3.---}
\emph{Let \(k\) be a field. The category of commutative affine
algebraic \(k\)-groups is abelian.}

Let us also point out that the category of all commutative affine
\(k\)-groups (not necessarily of finite type) is abelian. This
follows from Expos\'e~VI\(_B\), \SGARefLink{thm:sga3-VIB-11.17}{11.17} and \SGARefLink{thm:sga3-VIB-11.18.2}{11.18.2}
(cf.\ \textup{[DG70]}, \S~III.3, \SGARefLink{prop:sga3-VIB-7.4}{7.4}); see also
Expos\'e~VII\(_B\), 2.4.2 for a proof using formal groups.
